\subsection{Historical context}

Hadwiger's conjecture places planar coloring in a broader minor-theoretic
setting. The \(K_4\)-minor-free case was proved independently by
Hadwiger~\citep{hadwiger1943} and Dirac~\citep{Dirac1952}.
Wagner's structure theorem~\citep{Wagner1937} makes the
four-colorability of graphs without a \(K_5\) minor equivalent to the
Four-Color Theorem, proved by Appel and Haken, with Koch contributing
to the reducibility part~\citep{AppelHaken1977,AppelHakenKoch1977}.
Robertson, Seymour, and Thomas~\citep{RobertsonSeymourThomas1993}
proved that every graph without a \(K_6\) minor is five-colorable,
by reducing the obstruction to a graph that becomes planar after
one vertex is removed. These results explain the conjecture's link
between coloring and the structure of minor-closed classes.

For general \(t\), the density bounds of
Kostochka~\citep{Kostochka1984} and Thomason~\citep{Thomason1984}
give \(O(t\sqrt{\log t})\)-colorability for graphs with no \(K_t\)
minor. Norin, Postle, and Song~\citep[Theorem~1.4]{NorinPostleSong2023}
broke this degeneracy barrier, obtaining
\(O(t(\log t)^\beta)\) colors for every fixed \(\beta>1/4\).
Delcourt and Postle~\citep[Theorem~1.5]{DelcourtPostle2025} reduced the
bound to \(O(t\log\log t)\). A recent preprint of Liu and
Luo~\citep[Theorem~1.1]{LiuLuo2026} further gives
\(O(t\log\log\log t)\)-colorability.

The restriction of Hadwiger's conjecture to graphs of independence
number at most two has an equivalent formulation: every such
graph on \(m\) vertices should have a
\(K_{\lceil m/2\rceil}\) minor. Plummer, Stiebitz, and
Toft~\citep[Theorem~3.1]{pst2003} proved equivalence between the
restricted conjecture and this universally quantified minor assertion.
This identifies the case
\(\alpha(G)\le2\) as a setting in which chromatic number can be replaced by an
explicit fraction of the order.

Duchet and Meyniel~\citep{duchet1982} proved the general lower bound
\(h(G)\ge |V(G)|/(2\alpha(G)-1)\), giving \(h(G)\ge m/3\)
when \(\alpha(G)=2\). Blasiak~\citep{Blasiak2007} obtained complete-minor models of order
\(\Omega(m^{4/5})\) using only singleton and edge branch sets.
Fox~\citep[Theorem~2]{fox2010} strengthened the \(m/3\) bound to
\[
             h(G)\ge \frac m3+c\,m^{4/5}(\log m)^{1/5}
\]
for an absolute constant \(c>0\). His Theorem~3 also yields a
connected matching of size \(\Omega(m^{4/5}(\log m)^{1/5})\), since
its alternative clique of the same order can be paired. These bounds
describe the quantitative background
to the fixed linear upper bounds in \Cref{thm:main,cor:hadwiger}.

F\"uredi, Gy\'arf\'as, and Simonyi~\citep{fgs2005}
proposed the sharp connected-matching conjecture: a graph with
independence number two and order \(4t-1\) has a connected matching
of size \(t\). They verified the assertion through \(t=17\),
and Chen and Deng~\citep{ChenDeng2025} extended this
range to \(t=22\). Cambie~\citep[Theorem~1.5]{cambie2021} proved that
Hadwiger's conjecture implies this statement. The graphs in
\Cref{thm:main} also disprove the sharp connected-matching conjecture:
delete at most three vertices to obtain order \(4t-1\); the
connected-matching number does not increase and remains below \(t\)
for sufficiently large orders. The remaining graph has independence
number exactly two, since a complete graph of order \(4t-1\) has a
connected matching of size \(2t-1\).

Substantial positive results also come from additional clique structure.
Chudnovsky and Seymour~\citep[Theorem~1.3]{ChudnovskySeymour2012}
proved that an \(m\)-vertex graph with independence number at most two
has a \(K_{\lceil m/2\rceil}\) minor whenever its clique number is
at least \(m/4\) for even \(m\), or at least \((m+3)/4\) for odd
\(m\). Their method packs induced three-vertex paths. Norin and
Seymour~\citep[Theorem~1.2]{NorinSeymour2026} subsequently combined this
structure with random pairing to obtain minors of order \(\chi(G)\)
and edge density at least \(\gamma-o(1)\), where
\(\gamma=0.986882\ldots\) is a fixed constant.
These results highlight the requirement in a complete-minor model:
every pair of branch sets must be joined, whereas high density permits
missing adjacencies.

A weaker conjecture asks only for some absolute \(c_0>0\) such that
\(\cm(G)\ge c_0|V(G)|\) for all sufficiently large graphs with
\(\alpha(G)\le2\)~\citep{fgs2005}.
Yip~\citep[Theorem~7]{yip2025} proved a dense variant: a matching of
linear size can have the proportion of its non-touching edge pairs
made arbitrarily small, with the linear coefficient depending on the
prescribed proportion. That statement allows such pairs to remain.
The fixed coefficient in \Cref{thm:main} leaves the weaker conjecture
for smaller positive coefficients unresolved; it does not assert
\(\cm(G)=o(|V(G)|)\).

A related parity-constrained problem has a different recent resolution.
An odd \(K_t\) minor consists of \(t\) disjoint trees, each properly
two-colored, with a same-color edge joining every pair of trees.
K\"uhn, Sauermann, Steiner, and Wigderson~\citep{KuhnSauermannSteinerWigderson2025}
disproved the odd Hadwiger conjecture using graphs of independence
number at most two. The construction in this paper bounds the order
of ordinary complete minors.

The separate companion manuscript~\citep[Theorem~1.1]{OpenAIColinDeVerdiereCounterexample2026}
constructs graphs with \(\chi(G)>\mu(G)+1\), where \(\mu\) is the
ordinary Colin de Verdi\`ere invariant. Since \(h(G)\le\mu(G)+1\),
its result also gives counterexamples to Hadwiger's inequality.
The argument here proves the connected-matching bound directly and
does not use that invariant.
